\documentclass[10pt,a4paper]{article}

\usepackage[margin=1.5cm]{geometry}
\usepackage[T1]{fontenc}
\usepackage[utf8]{inputenc}
\usepackage{mathptmx} % Times-like font
\usepackage{enumitem}
\usepackage{xcolor}
\usepackage[hidelinks]{hyperref}

\pagestyle{empty}
\setlength{\parindent}{0pt}
\setlist[itemize]{leftmargin=0.8em, labelindent=0pt, labelwidth=0.5em, labelsep=0.3em, align=left,
  itemsep=0pt, topsep=1pt, parsep=0pt, label=\textbullet}

% ---------- Layout macros ----------
\definecolor{rulegray}{gray}{0.65}
\newlength{\cvlabelwidth}\setlength{\cvlabelwidth}{3cm}

% Left-column label; zero depth so multi-line labels don't push the content down
\newcommand{\cvlabel}[1]{%
  \raisebox{0pt}[\height][0pt]{\parbox[t]{\cvlabelwidth}{\raggedright\textbf{#1}}}\hfill}

% \cvsection{LABEL}{content}: label in the left column, content on the right
\newcommand{\cvsection}[2]{%
  \par\noindent\textcolor{rulegray}{\rule{\linewidth}{0.4pt}}\par\vspace{5pt}
  \begin{list}{}{\setlength{\leftmargin}{\cvlabelwidth}\setlength{\labelwidth}{\cvlabelwidth}%
      \setlength{\labelsep}{0pt}\setlength{\topsep}{0pt}\setlength{\itemsep}{0pt}%
      \setlength{\parsep}{0pt}}
    \item[{\cvlabel{#1}}]#2
  \end{list}
  \vspace{8pt}
}

% \entry{title}{date}{role}{place}: two lines, left and right aligned
\newcommand{\entry}[4]{%
  \textbf{#1}\hfill\textbf{#2}\par
  \ifx&#3&\else\textit{#3}\hfill\textit{#4}\par\fi
}

% \project{name}{url}{description}
\newcommand{\project}[3]{%
  \textbf{#1}\quad\href{#2}{\underline{#1}}\par
  \begin{itemize}\item #3\end{itemize}
  \vspace{3pt}
}

\newcommand{\gap}{\vspace{5pt}}

\begin{document}

% ---------- Header ----------
{\raggedleft\fontsize{24}{28}\selectfont JOÃO ALVIM\par}
\vspace{4pt}
{\small\noindent
Braga, Portugal \textbar{} +351 934662832 \textbar{}
\href{mailto:joaoalvim179@proton.me}{\underline{joaoalvim179@proton.me}} \textbar{}
\href{https://www.linkedin.com/in/jo\%C3\%A3o-alvim-050778223}{\underline{linkedin/joao-alvim}} \textbar{}
\href{https://github.com/jotaalvim}{\underline{github.com/jotaalvim}} \textbar{}
\href{https://jotaalvim.github.io}{\underline{jotaalvim.github.io}}\par}
\vspace{6pt}

% ---------- Education ----------
\cvsection{EDUCATION}{%
  \entry{University of Minho}{October 2020 - October 2025}{Bachelors and Masters in Computer Engineering, thesis in Formal Methods}{Portugal}\gap
  Thesis: \textit{[Abstract Debugging for WebAssembly]} \par\gap
  \entry{Ghent University, Belgium}{October 2024 - March 2025}{Masters - Erasmus+ (1 semester)}{Belgium}\gap
%  \entry{Escola Secundária Alberto Sampaio}{September 2017 - June 2020}{Secondary Education, Sciences}{}\gap
}

% ---------- Work experience ----------
\cvsection{WORK\\EXPERIENCE}{%
  \entry{INVISIBLE LAB}{January 2026 - Present}{Backend Software Engineer}{}
  Developed robust backend services in Golang backed by PostgreSQL. Designed external API
  integrations and delivered key features for Infocarro and a non-profit association platform.\par\gap

  \entry{DESY}{June 2024 - August 2024}{Summer Internship}{Hamburg, Germany}
  Developed a monitoring dashboard with Grafana for ASAPO -- high performance distributed streaming platform for DESY a German research center operating particle accelerators.
  \par\gap

  \entry{Programming Tutoring to Undergraduates}{2022 - 2025}{}{Braga}
  Taught Haskell and C to 4 and 5 students, Natural Language Processing with Python to 1 student,
  and statistics with R.\par\gap

}

% ---------- Publications ----------
\cvsection{PUBLICATIONS}{%
  \begin{itemize}[topsep=0pt]
    \item Online Abstract Debugging - Out of the loop no more. SLATE 2026 -- A debugging framework that allows abstract
      interpretation to be applied on-demand, gaining precision at specific points of program
      analysis before obtaining the full static results. It allows the user to unroll loops and
      explore intermediate states where one would normally lose precision. \quad
      \href{https://doi.org/10.4230/OASIcs.SLATE.2026.20}{\underline{https://doi.org/10.4230/OASIcs.SLATE.2026.20}}
  \end{itemize}
}

% ---------- Projects ----------
\cvsection{PROJECTS}{%
  \project{pocket-cube-solver}{https://github.com/jotaalvim/pocket-cube-solver}{Haskell solution to the Pocket Cube from Justin.}

  \project{cbirds}{https://github.com/jotaalvim/cbirds}{Python module demonstrating combinator birds in point-free style.}

  \project{QForms}{https://github.com/jotaalvim/QForms}{QuickForms is a self-hosted local form generation tool designed as an alternative to proprietary services like Google Forms. I've used Flask and my Python module is on PyPI.}

  \project{vscode-ipma}{https://github.com/jotaalvim/ipma-vsc/tree/main/ipma-carta\#readme}{VS Code extension that shows a popup with a beautiful IPMA's weather map.}
  \project{shipslogs}{https://github.com/jotaalvim/shipslogs}{Python tool that simplifies keeping records of screenshots taken during events such as talks, meetings or online classes.}
  \project{Chess-Heatmaps}{https://github.com/jotaalvim/chess-heatmaps}{Represent chess moves as heat maps in Python.}

  %\project{Instagram Puzzles Generator}{https://github.com/jotaalvim/instagram-puzzles-generator}{Automates generation of chess puzzle images published on my university chess club's Instagram.}
  \project{chessredditstories}{https://www.tiktok.com/@chessredditstories}{Automated content creation that generates reels combining chess gameplay with Reddit stories using Python.}
  %\project{datalist}{https://github.com/jotaalvim/datalist}{Implementation of all functions in Haskell's Data.List module.}
  %\project{hmusik}{https://github.com/jotaalvim/hmusik}{Generating chord progressions with Haskell.}
}

% ---------- Technical skills ----------
\cvsection{TECHNICAL\\SKILLS}{%
  \textbf{Languages}: Portuguese (fluent), English (B2 exam)\par
  \textbf{Programming}: Go, PostgreSQL, Haskell, C, C++, Python, Java, TypeScript, JavaScript, LaTeX, SQL, Linux, VIM,
  WebAssembly, Machine Learning
}

% ---------- Awards ----------
\cvsection{AWARDS}{%
  \begin{itemize}[topsep=0pt]
    \item \textbf{2nd place:} SEI 2021 Hackathon (shipslogs)
    \item \textbf{1st place:} SEI 2025 \href{https://drive.google.com/drive/folders/1xnQhL0wbuDtT2qomy4DPXgj36mWrUCKh}{\underline{Programming Competition}}
  \end{itemize}
}

% ---------- Volunteer ----------
\cvsection{VOLUNTEER}{%
  \entry{Camerata Minhota}{January 2023 - 2025}{}{}
  Co-founded a chamber music group at the University of Minho, performing as pianist at university and public events.\par\gap
  \entry{Coderdojo}{January 2020 - January 2023}{Mentored children learning to code}{}\gap
  \entry{Recriativo, Núcleo de Estudantes de Engenharia Informática}{2021 - 2022}{Collaborator}{}
}


% ---------- Extracurricular ----------
\cvsection{EXTRA-\\CURRICULAR}{%
  \begin{itemize}[topsep=0pt]
    \item Conservatório Bomfim - September 2017 - June 2020 - Secondary Education, Music -- 8th Grade in Piano
    \item Led a cat rescue from a 15 meter tree:
      \href{https://ominho.pt/braga-escalador-resgata-gato-preso-ha-seis-dias-em-arvore-com-15-metros/}{\underline{Climber rescues cat trapped in 15 meter tree for six days}}

    \item Attended \href{https://ifl23.github.io/}{\underline{IFL 2023}} -- Implementation and Application of Functional Languages
        conference, SLATE 2026

    \item Federated chess and climbing athlete

    \item Chess Coach - 2022 - 2024 - University of Minho Chess Club \& Tempos Brilhantes: Taught weekly lessons analyzing classic positions and organized tournaments for the university chess club; taught biweekly extracurricular chess classes to primary school students.\par\gap

    \item Seasonal Outdoor Work - 3 weeks - High-Altitude Work Technician; Monitor \& Diver at Lanhoso Adventure Park\gap



  \end{itemize}
}
\end{document}
